\section{Conclusion}

\subsection{What was achieved}

The project delivers a complete, working compiler front end for a language that has no
standard written form:

\begin{itemize}[leftmargin=1.4em, itemsep=0.15em]
  \item A corpus of \CorpusTotal\ statements heard in Yaound\'e, transcribed by hand,
        cross-checked, and attributed across all three team members, with the attestation
        enforced in software rather than asserted in prose.
  \item A hand-written scanner and lexer built on five documented regular expressions,
        resolving \LexiconTotal\ lexicon entries across \OriginCount\ donor languages into
        \TerminalCount\ terminal classes, with multiword recognition, accent folding and
        homograph protection.
  \item A context-free grammar of \ProductionCount\ productions authored in natural
        left-recursive form, mechanically transformed in \LedgerSteps\ recorded steps into
        an \ExecProductionCount-production LL(1) grammar.
  \item FIRST and FOLLOW sets and a parsing table verified conflict-free by an automated
        test.
  \item A table-driven predictive parser returning an explicit derivation on acceptance and
        an expected-terminal set plus concrete repair suggestions on rejection, accepting
        \CorpusAccepted\ of \CorpusTotal\ collected statements.
  \item A web application exposing every one of those artefacts, so the compiler is
        demonstrable rather than merely described.
\end{itemize}

\subsection{What we learned}

The most useful lesson was that the transformations taught in the course are not
bookkeeping. Eliminating left recursion on \nonterm{Seq} \emph{created} a new common prefix
on \term{SEP} that had to be factored out afterwards --- the two techniques interact, and the
order in which they are applied decides whether the result is LL(1). Seeing that happen on
our own grammar, in the ledger the tool recorded, taught more than the worked examples did.

The second lesson concerns honesty in tooling. It was tempting, each time a statement was
rejected, to add a production until it passed. We resisted, and the two remaining
rejections in Section~\ref{sec:results} are documented with the specific LL(1) conflict
that prevents a fix. A grammar that accepts everything distinguishes nothing.

The third was methodological: generating every table in this document from the running
system, rather than transcribing results by hand, caught several places where earlier
documentation had drifted away from the code it claimed to describe.

\subsection{Limitations}

\begin{itemize}[leftmargin=1.4em, itemsep=0.15em]
  \item Two constructions --- post-object adverbs and post-verbal clitics --- lie outside the
        LL(1) grammar, for the reasons given in Section~\ref{sec:results}.
  \item The lexicon's dictionary tier (\LexiconDictionary\ of \LexiconTotal\ entries) is
        reference-derived rather than corpus-attested, and is labelled as such.
  \item Analysis is purely syntactic. The analyzer recognises structure; it does not
        interpret meaning.
\end{itemize}

\subsection{Future work}

Splitting \term{PRON} into full pronouns and clitics would remove both documented
rejections at the cost of re-annotating the pronoun entries against fresh evidence. Beyond
that, the natural next step is a semantic layer: the lexicon already carries glosses and
donor languages, so a gloss-composition pass over the parse tree would turn the analyzer
from a recogniser into a translator --- which is what would make it useful outside a
coursework setting. A larger corpus would also let the topic matcher be evaluated for
precision and recall rather than merely cited for evidence.

\begin{todobox}[Outstanding before submission]
\begin{enumerate}[leftmargin=1.4em, itemsep=0.15em]
  \item \textbf{Presentation.} Prepare the ten-minute slide deck and live demonstration
        (Component 5C).
  \item \textbf{Deployment hardening.} Replace the placeholder registration code before the
        public instance is exposed.
\end{enumerate}
\end{todobox}
